%% ---------------------------------------------------------------------------
%% apendice.tex
%%
%% Apéndice A: implicaciones ambientales y de desarrollo sostenible (obligatorio
%% según los Lineamientos V4.1; estructura del TFG de ejemplo). Apéndices B y C:
%% sin contenido.
%%
%% Convención de escritura: una oración por línea del archivo fuente.
%% ---------------------------------------------------------------------------

\chapter{APENDICE 1: Implicaciones ambientales y de desarrollo sostenible}
\label{appendix:apendiceA}

Este apéndice analiza las implicaciones de la plataforma SLAM híbrida en tiempo real sobre la sociedad y sobre el ejercicio profesional de la ingeniería en computadores.
Para cada aspecto se identifican los impactos actuales y potenciales, se relacionan con los Objetivos de Desarrollo Sostenible (ODS) de las Naciones Unidas \cite{onu2015agenda2030} y se plantean proyecciones para el área de la localización y la fusión sensorial.
El análisis parte de la plataforma desarrollada y de sus resultados (Capítulo~\ref{ch:resultados}), no de un sistema ideal.

\section{Aspectos sociales}

\textbf{Identificación y análisis:}
La plataforma estima la posición de un vehículo en tiempo real con los sensores de un teléfono inteligente de consumo y una computadora con tarjeta gráfica, sin equipos especializados de posicionamiento, lo que reduce la barrera de entrada para experimentar con localización y fusión sensorial en ámbitos académicos.
Además, queda documentada como una plataforma experimental para la Escuela de Ingeniería en Computadores.
El principal riesgo social es que se confíe en la estimación más de lo que corresponde: sin GPS, el error crece hasta cientos o miles de metros en pocos minutos, en los cortes de 30~s la mediana del error al volver el GPS llega a 101~m en la configuración completa (Tabla~\ref{tab:res-modos}), y la política de aprendizaje por refuerzo no se traslada a un dominio con otra aplicación y otra cámara.
Por ello, la plataforma no debe usarse como única fuente de posición para tomar decisiones de conducción.

\textbf{Relación con los ODS:}
se relaciona con el ODS 4 (educación de calidad), por su uso como herramienta de experimentación e investigación, y con el ODS 11 (ciudades y comunidades sostenibles), por su aporte a la localización en aplicaciones de movilidad.

\textbf{Proyecciones:}
con más recorridos y una referencia de mayor precisión, la plataforma podría apoyar la investigación en asistencia a la conducción y robótica móvil con sensores de bajo costo.
Para ello debe comunicar la incertidumbre de su estimación, de modo que quien la use sepa cuándo no confiar en ella.

\section{Aspectos económicos}

\textbf{Identificación y análisis:}
La plataforma se construyó con herramientas gratuitas o de código abierto, como Python, OpenCV, PyTorch y Android Studio, y con el modelo de profundidad Depth Anything~V2 \cite{yang2024depthanything}, y se ejecuta sin conexión a internet (requerimiento RC-1), de modo que no depende de servicios de pago en la nube.
Los sensores son los de un teléfono de consumo, en lugar de receptores de posicionamiento de alta precisión o de sensores LiDAR.
El costo principal es el cómputo: la plataforma requiere una tarjeta gráfica NVIDIA, y la portátil empleada en campo, de menor capacidad, procesa cada cuadro de dos a tres veces más lento que la computadora de escritorio, con una latencia mayor.

\textbf{Relación con los ODS:}
se relaciona con el ODS 9 (industria, innovación e infraestructura), por el uso de tecnología de bajo costo y de código abierto para la innovación en localización.

\textbf{Proyecciones:}
un uso a mayor escala tendría que valorar el costo de la tarjeta gráfica frente al de sensores especializados.
Ejecutar la plataforma en equipos de menor costo exigiría optimizar la odometría visual, que es el cuello de botella de la latencia en la portátil.

\section{Aspectos legales}

\textbf{Identificación y análisis:}
La plataforma usa datos y modelos de terceros con condiciones propias.
El conjunto KITTI \cite{geiger2012kitti}, empleado en los estudios de la segunda fase, se publica bajo la licencia Creative Commons Atribución-NoComercial-CompartirIgual 3.0 y para uso académico \cite{kitti_web}, por lo que los resultados obtenidos con él no pueden explotarse comercialmente.
El modelo Depth Anything~V2 se distribuye bajo la licencia Apache~2.0 en su variante pequeña y bajo una licencia no comercial en las de mayor tamaño \cite{depthanything_repo}; la plataforma usa una variante pequeña ajustada para profundidad métrica en exteriores, cuya licencia debe confirmarse antes de cualquier uso comercial.
Las grabaciones propias contienen las posiciones GPS de los recorridos y video de vías públicas en el que pueden aparecer personas y placas de vehículos, y en Costa Rica cualquier dato relativo a una persona física identificada o identificable es un dato personal \cite{ley8968}.

\textbf{Relación con los ODS:}
se relaciona de forma indirecta con el ODS 16 (paz, justicia e instituciones sólidas), por el respeto a la propiedad intelectual y a la privacidad en el manejo de los datos.

\textbf{Proyecciones:}
una publicación abierta de la plataforma, de sus grabaciones o de sus resultados exigirá anonimizar los rostros y las placas en el video, valorar si las trayectorias revelan datos personales, como lugares de residencia o de trabajo, y respetar la licencia no comercial de KITTI.

\section{Aspectos ambientales}

\textbf{Identificación y análisis:}
La plataforma reutiliza un teléfono y computadoras de uso general en lugar de equipo especializado, y procesa todo localmente, sin servidores de inferencia en funcionamiento permanente.
Además, cada grabación se procesa una sola vez en el modo determinista y se guarda en una caché, de modo que las evaluaciones posteriores vuelven a ejecutar la fusión sin procesar de nuevo el video.
El principal impacto negativo es el consumo de energía del cómputo con tarjeta gráfica: el modelo de profundidad se ejecuta en la GPU durante toda la operación, y el entrenamiento de cada política de aprendizaje por refuerzo tarda unos 12~min, repetido para las tres semillas y las siete pruebas de cada iteración; este consumo no se midió.
Las pruebas en vehículo implican, además, el consumo de combustible de los recorridos.

\textbf{Relación con los ODS:}
se relaciona con el ODS 12 (producción y consumo responsables), por la reutilización de dispositivos de consumo en lugar de equipo dedicado, y con el ODS 13 (acción por el clima), por el consumo energético del cómputo.

\textbf{Proyecciones:}
medir el consumo de energía y la utilización de recursos permitiría valorar el costo ambiental de cada configuración.
Un modelo de profundidad más liviano reduciría ese consumo, a cambio de una posible pérdida de precisión en la escala.

\section{Aspectos de seguridad}

\textbf{Identificación y análisis:}
Desde la perspectiva de la seguridad de la información, el protocolo entre el teléfono y la computadora no incluye autenticación ni cifrado: el teléfono actúa como servidor, atiende a una sola computadora a la vez y una conexión nueva reemplaza a la anterior, de modo que cualquier equipo en la misma red podría recibir el video y las posiciones o interrumpir la sesión.
En las pruebas de campo, la conexión por cable mediante anclaje de red por USB limitó ese acceso a la computadora conectada.
Desde la perspectiva de la seguridad física, las pruebas se realizan en vías públicas, con el teléfono montado en el tablero y una computadora dentro del vehículo, lo que exige que el soporte no obstruya la visibilidad y que el conductor no opere la computadora mientras conduce.
Además, los errores de la estimación sin GPS, de cientos de metros, impiden usarla para guiar un vehículo.

\textbf{Relación con los ODS:}
se relaciona con el ODS 9 (industria, innovación e infraestructura), por la confiabilidad de los sistemas de localización, y con el ODS 11 (ciudades y comunidades sostenibles), en su meta de sistemas de transporte seguros.

\textbf{Proyecciones:}
un despliegue fuera del entorno experimental requerirá autenticar la conexión entre el teléfono y la computadora, cifrar los datos y advertir de forma explícita sobre la incertidumbre de la estimación durante las pérdidas del GPS.

\section{Aspectos de salud}

\textbf{Identificación y análisis:}
La plataforma no es un dispositivo médico ni procesa información de salud.
Sus efectos sobre la salud son indirectos y se relacionan con las pruebas en vehículo: la distracción del conductor puede provocar accidentes, por lo que la operación de la computadora y la revisión de la interfaz deben quedar a cargo de un acompañante o realizarse con el vehículo detenido.
A largo plazo, una localización más estable podría contribuir a sistemas de asistencia a la conducción que reduzcan los accidentes, aunque la plataforma actual no tiene ese grado de madurez.

\textbf{Relación con los ODS:}
se relaciona de forma limitada con el ODS 3 (salud y bienestar), en su meta de reducir las muertes y lesiones causadas por accidentes de tránsito, a la que solo podría aportar una versión madura de la plataforma.

\textbf{Proyecciones:}
una versión futura debería acompañarse de un procedimiento de pruebas en vía con funciones definidas para el conductor y el operador, y de una interfaz que no requiera atención durante la conducción.

\chapter{APENDICE B}
\label{appendix:apendiceB}



\chapter{APENDICE C}
\label{appendix:apendiceC}

%%% Local Variables:
%%% mode: latex
%%% TeX-master: "main"
%%% End:
